\chapter{1000 तक की संख्याएँ}\label{ch:g2:numbers}

पिछले साल हमने इकाइयों को दहाइयों में बाँधा था (\cref{ch:g1:tens});
अब दस दहाइयाँ मिलकर एक \emph{सैकड़ा} बनती हैं, और संख्याएँ $1000$
तक पहुँच जाती हैं। तीन अंक, तीन काम: सैकड़ा, दहाई, इकाई।

\section{सैकड़ा, दहाई, इकाई}

\begin{definition}[तीन स्थान]\label{def:g2:numbers:places}
एक \emph{सैकड़ा} दस दहाइयाँ हैं, यानी दस के दस बंडल। संख्या
$346$ के तीन हिस्से हैं:
\[
346 = 3 \text{ सैकड़ा} + 4 \text{ दहाई} + 6 \text{ इकाई}
= 300 + 40 + 6 .
\]
\end{definition}

\begin{omfigure}
\begin{tikzpicture}[scale=0.85]
  \draw (0,0) grid[xstep=2.3, ystep=0.85] (6.9,1.7);
  \node[font=\small] at (1.15,1.275) {सैकड़ा};
  \node[font=\small] at (3.45,1.275) {दहाई};
  \node[font=\small] at (5.75,1.275) {इकाई};
  \node[font=\large, omDef] at (1.15,0.425) {$3$};
  \node[font=\large, omDef] at (3.45,0.425) {$4$};
  \node[font=\large, omDef] at (5.75,0.425) {$6$};
\end{tikzpicture}

{\small $346$ की स्थान-सारणी; इसे ``तीन सौ छियालीस'' पढ़ते हैं।
(इससे बड़ी संख्याएँ \cref{ch:g3:numbers} में मिलेंगी।)}
\end{omfigure}

\begin{example}[शून्य स्थान रोके रखता है]\label{ex:g2:numbers:zero}
``तीन सौ पाँच'' यानी $305$: तीन सैकड़ा, \emph{कोई} दहाई नहीं,
पाँच इकाई। $0$ ही $5$ को इकाई के स्थान पर रोके रखता है। उसके
बिना $35$ बिलकुल दूसरी संख्या हो जाती।
\end{example}

\begin{example}[अदला-बदली]\label{ex:g2:numbers:exchange}
दस इकाई मिलकर एक दहाई बनती हैं, दस दहाई मिलकर एक सैकड़ा। तो
$10$ दहाई $= 100$, और $2$ सैकड़ा $= 20$ दहाई। संख्या $250$ में
$25$ दहाइयाँ हैं, केवल अंक $5$ नहीं।
\end{example}

\section{तुलना और क्रम}

\begin{method}[तीन अंकों वाली संख्याओं की तुलना]\label{met:g2:numbers:compare}
\begin{enumerate}
  \item जिसमें ज़्यादा अंक, वही बड़ी: $103 > 98$;
  \item अंक बराबर हों: बाईं ओर से तुलना करो --- पहले सैकड़ा, फिर
        दहाई, फिर इकाई। पहला \omterm{ex:g1:subtraction:difference}{अंतर} ही फ़ैसला कर देता है:
        $472 > 468$, क्योंकि दहाई पर $7 > 6$।
\end{enumerate}
\end{method}

\begin{example}\label{ex:g2:numbers:order}
$530$, $305$, $503$ को क्रम में लगाओ: पहले सैकड़ा ($5$, $3$, $5$),
इसलिए $305$ सबसे आगे; $530$ और $503$ में दहाई फ़ैसला करती है ($3 > 0$):
\[
305 < 503 < 530 .
\]
\end{example}

\section{संख्या रेखा पर छलाँगें}

\begin{method}[दहाई और सैकड़ा जोड़ना]\label{met:g2:numbers:jumps}
रेखा पर:
\begin{enumerate}
  \item $+1$ / $-1$: इकाई बदलती है;
  \item $+10$ / $-10$: दहाई बदलती है ($362 + 10 = 372$);
  \item $+100$ / $-100$: सैकड़ा बदलता है ($362 + 100 = 462$)।
\end{enumerate}
हर बार केवल एक अंक हिलता है --- सिवाय तब, जब कोई स्थान ``भर
जाए'' और हासिल आगे चला जाए, जैसे $195 + 10 = 205$।
\end{method}

\begin{omfigure}
\begin{tikzpicture}[scale=1.0]
  \draw[-Stealth] (-0.3,0) -- (10.6,0);
  \foreach \x/\l in {0.5/{362}, 5.5/{372}, 9.5/{472}}
    { \draw (\x,-0.1) -- (\x,0.1);
      \node[below, font=\small] at (\x,-0.12) {$\l$}; }
  \draw[-Stealth, omDef, thick] (0.5,0.28) .. controls (2.5,1.0) ..
    (5.5,0.28);
  \node[omDef, font=\small] at (3.0,1.05) {$+10$};
  \draw[-Stealth, omThm, thick] (5.5,0.28) .. controls (7.5,1.05) ..
    (9.5,0.28);
  \node[omThm, font=\small] at (7.5,1.1) {$+100$};
\end{tikzpicture}

{\small $362$ से छलाँगें: $+10$ की छोटी छलाँग दहाई बदलती है,
$+100$ की बड़ी छलाँग सैकड़ा।}
\end{omfigure}

\section{अभ्यास}

\begin{exercise}[$\star$]\label{exo:g2:numbers:1}
कितने सैकड़ा, दहाई और इकाई: $258$; \ $470$; \ $600$; \
$909$?
\end{exercise}

\begin{exercise}[$\star$]\label{exo:g2:numbers:2}
संख्या लिखो: $4$ सैकड़ा, $0$ दहाई, $7$ इकाई; \ $6$ सैकड़ा
और $3$ दहाई; \ दस सैकड़ा।
\end{exercise}

\begin{exercise}[$\star$]\label{exo:g2:numbers:3}
अंकों में लिखो: दो सौ सोलह; \ पाँच सौ चार; \
सात सौ नब्बे।
\end{exercise}

\begin{exercise}[$\star$]\label{exo:g2:numbers:4}
\cref{def:g2:numbers:places} की तरह हिस्सों में बाँटो: $438$; \ $702$; \
$560$।
\end{exercise}

\begin{exercise}[$\star$]\label{exo:g2:numbers:5}
उतार कर $<$ या $>$ से पूरा करो:
\[
327 \;?\; 372, \qquad
580 \;?\; 508, \qquad
199 \;?\; 200, \qquad
640 \;?\; 604 .
\]
\end{exercise}

\begin{exercise}[$\star$]\label{exo:g2:numbers:6}
सबसे छोटी से सबसे बड़ी के क्रम में लिखो: $430$; \ $304$; \ $340$; \ $403$।
\end{exercise}

\begin{exercise}[$\star$]\label{exo:g2:numbers:7}
रेखा पर छलाँग लगाकर निकालो: $253 + 10$; \quad $253 + 100$;
\quad $480 - 10$; \quad $716 - 100$।
\end{exercise}

\begin{exercise}[$\star$]\label{exo:g2:numbers:8}
आगे गिनो: $195, 205, 215, ?, ?$। \quad
$300, 400, 500, ?, ?$। \quad पहली वाली में सावधान!
\end{exercise}

\begin{exercise}[$\star$]\label{exo:g2:numbers:9}
एक दुकान में $100$-$100$ पेंसिलों के $6$ डिब्बे और $5$ खुली
पेंसिलें हैं। कुल कितनी पेंसिलें?
\end{exercise}

\begin{exercise}[$\star\star$]\label{exo:g2:numbers:10}
अंकों $2$, $5$ और $8$ को एक-एक बार लेकर तीन अंकों वाली सबसे
बड़ी और सबसे छोटी संख्या लिखो। दोनों में कितना फ़ासला है?
(चाहो तो रेखा पर छलाँग लगाकर निकाल सकते हो।)
\end{exercise}
